\documentclass{article}
\usepackage[utf8]{inputenc}
\usepackage{geometry}
\geometry{a4paper, margin=1in}

\title{The Shared Future Narrative Hypothesis}
\author{Albert Jan van Hoek}
\date{\today}

\begin{document}

\maketitle

I would like to share a hypothesis.

In my \textbf{Theory of Long-term Collaboration (TLC)}, I treat collaboration as an interaction between two learning networks — human brains.

Last week I learned that neuroscientists believe the brain learns through predictive coding: your brain constantly makes predictions about reality and updates itself when those predictions are wrong. In a collaboration, this means both people are continuously predicting each other's behavior — and then act on those predictions.

\section*{The Painful Trap of Predictions}

This can create a painful trap. When we expect the other to judge us negatively, we may already start acting as if that prediction is true. Emotions like shame, guilt, fear, embarrassment, resentment, and grief are not just feelings — under this framework they are predictions about how the other sees us, or how the future will unfold. And those predictions cause real present pain. The broken relationship hurts not just because of what happened, but because of what we expect to happen next.

\section*{The Shared Future Narrative}

This is the core of what I call the \textbf{Shared Future Narrative hypothesis}. To improve a collaboration — or make a broken one possible again — both people should address that negative prediction directly, by speaking about a happy future collaboration together. When the predicted future becomes a happy collaboration, deviating from that becomes a prediction error. Both parties then act to minimize that error — and the future becomes a self-fulfilling prophecy. The shared narrative dissolves the present pain, before any formal reconciliation has even taken place.

It follows the same logic as the golden rule: treat your neighbour as you want to be treated. In a system of collaborative predictive coders, it might be the mathematically stable solution.

\section*{Application to Partisan Divides}

So how to resolve a partisan divide between Democrats and Republicans? Talk about it as if the healing is already underway. ``Do you remember 2026, when we stopped predicting the worst of each other and started building something together again?'' Such narratives should begin to rewrite what both systems expect — and therefore what both sides do.

\section*{A Call to Action}

For now this is still a hypothesis. Perhaps worth testing though. Let's organize a reunion-party?

\end{document}
